\documentclass[aspectratio=169]{beamer}

\usepackage[T1,T2A]{fontenc}
\usepackage[utf8]{inputenc}
\usepackage[main=russian,english]{babel}
\usepackage{cmap}
\usepackage{tikz}
\usetikzlibrary{arrows.meta,positioning}
\usepackage{array}
\usepackage{booktabs}
\usepackage{amssymb,amsfonts,amsmath}
\usepackage{listings}

\usecolortheme{beaver}
\usefonttheme[onlymath]{serif}
\setbeamertemplate{navigation symbols}{}
\setbeamertemplate{footline}[page number]
\setbeamertemplate{blocks}[rounded=true,shadow=true]
\setbeamersize{text margin left=0.5cm,text margin right=0.5cm}
\setbeamerfont{author}{size=\normalsize}
\setbeamerfont{institute}{size=\normalsize}
\setbeamerfont{date}{size=\normalsize}
\setbeamertemplate{enumerate item}{\insertenumlabel)}
\setbeameroption{hide notes}

\newcommand{\E}{\mathbb E}
\newcommand{\sg}{\operatorname{stopgrad}}
\newcommand{\KL}{\operatorname{KL}}
\newcommand{\codefile}[1]{\par\medskip{\scriptsize\color{gray}\texttt{\detokenize{#1}}\par}}
\lstset{language=Python,basicstyle=\ttfamily\footnotesize,
  keywordstyle=\color{darkred},commentstyle=\color{gray},
  showstringspaces=false,columns=fullflexible,keepspaces=true,
  breaklines=true,frame=none}

\title{ReMix: reparameterization for mixtures}
\subtitle{Proof of Concept}
\author{Команда проекта: Александр Уденеев, Алексей Пономарев,\\Лина Горбунова, Максим Иванов}
\institute{ИАД, МФТИ}
\date{\empty}

\begin{document}

\begin{frame}
  \thispagestyle{empty}
  \maketitle
\end{frame}

\begin{frame}{Этап прототипа}
  \begin{block}{Цель этапа}
    Реализовать один простой оценщик в интерфейсе будущей библиотеки
    и подготовить небольшой воспроизводимый эксперимент.
  \end{block}
  \begin{enumerate}
    \item Представить смесь независимо от семейства компонент.
    \item Получить скалярную функцию потерь для \texttt{backward()}.
    \item Сохранить общий интерфейс для следующих оценщиков.
  \end{enumerate}
  \medskip
  Реализован ScoreFunction как простой пример оценщика
\end{frame}

\begin{frame}{Повторение про область задачи}
  \[
    q_\theta(z)=\sum_{k=1}^{K}\pi_k q_{k,\phi_k}(z),
    \qquad \pi=\operatorname{softmax}(a),\quad
    \theta=(a,\phi_1,\ldots,\phi_K).
  \]
  
  \begin{block}{Условия совместимости}
    \begin{enumerate}
        \item Одинаковый shape события.
        \item Общая базовая мера.
        \item Согласованные устройства и типы.
        \item Компоненты без batch-измерений.
    \end{enumerate}
  \end{block}

  \medskip
  Оценщик принимает \texttt{Mixture}~--- отдельное распределение отклоняется.
\end{frame}

\begin{frame}{Архитектура текущей реализации}
  \centering
  \begin{tikzpicture}[
    box/.style={draw=darkred,rounded corners,align=center,
      text width=3.65cm,minimum height=1cm,font=\small},
    link/.style={-{Latex[length=2mm]},thick,draw=darkred}]
    \node[box] (q) {\texttt{Mixture}\\\texttt{sample(), log\_prob()}};
    \node[box,right=0.65cm of q] (est) {\texttt{ScoreFunction}\\\texttt{estimate(q, objective)}};
    \node[box,right=0.65cm of est] (result) {\texttt{EstimateResult}\\\texttt{value, surrogate\_loss}};
    \node[box,below=0.65cm of q] (params) {Параметры пользователя:\\логиты и параметры компонент};
    \node[box,below=0.65cm of est] (obj) {\texttt{ELBOObjective}\\\texttt{evaluate(q, z)}};
    \node[box,below=0.65cm of result] (opt) {\texttt{backward()}\\\texttt{torch.optim.Adam}};
    \draw[link] (params) -- (q);
    \draw[link] (q) -- (est);
    \draw[link] (obj) -- (est);
    \draw[link] (est) -- (result);
    \draw[link] (result) -- (opt);
  \end{tikzpicture}
  \par\medskip\raggedright
  После шага оптимизатора заново создаётся смесь
  с текущими параметрами.
  \codefile{src/remix/}
\end{frame}

\begin{frame}{Как представлена смесь}
  \texttt{Mixture(logits, components)} хранит список распределений PyTorch.
  \[
    \ell_k(z)=\log\pi_k+\log q_k(z),\qquad
    \log q(z)=\operatorname{logsumexp}_k\ell_k(z).
  \]
  \[
    r_k(z)=\operatorname{softmax}_k\ell(z)
           =\frac{\pi_kq_k(z)}{q(z)}.
  \]
  \begin{itemize}
    \item \texttt{log\_prob}: устойчивая агрегация плотностей компонент.
    \item \texttt{responsibilities}: вероятности принадлежности компонентам.
    \item \texttt{sample}: выбор индекса из Categorical и значения компоненты.
  \end{itemize}
  \medskip
  Генерируются кандидаты от всех компонент, затем выбирается один.
  \codefile{src/remix/distributions.py}
\end{frame}

\begin{frame}{Целевая функция}
  \[
    \mathcal L(\theta)=\E_{z\sim q_\theta}[f_\theta(z)],
    \qquad f_\theta(z)=\log p(x,z)-\log q_\theta(z).
  \]
  \begin{block}{\texttt{ELBOObjective(log\_joint)}}
    Принимает пользовательскую функцию $\log p(x,z)$.
    Метод \texttt{evaluate(q, z)} возвращает значения $f_\theta(z)$.
  \end{block}
  \begin{center}
    \begin{tabular}{ll}
      \toprule
      Поле результата & Назначение \\
      \midrule
      \texttt{value} & Среднее $f_\theta(z)$ без градиента \\
      \texttt{surrogate\_loss} & Скаляр для минимизации через \texttt{backward()} \\
      \bottomrule
    \end{tabular}
  \end{center}
  \medskip
  В бенчмарке $p$ --- нормированное целевое распределение:
  $\mathcal L(\theta)=-\KL(q_\theta\Vert p)$.
  \codefile{src/remix/objectives.py}
\end{frame}

\begin{frame}{Score-function}
  При допустимости дифференцирования под интегралом:
  \[
    \nabla_\theta\mathcal L
      =\E_{q_\theta}\!\left[
        f_\theta(z)\nabla_\theta\log q_\theta(z)
        +\partial_\theta f_\theta(z)\right].
  \]
  \[
    \widehat g=\frac1N\sum_{i=1}^N
      \left[(f_\theta(z_i)-b)\nabla_\theta\log q_\theta(z_i)
        +\partial_\theta f_\theta(z_i)\right],
    \quad z_i\sim q_\theta.
  \]
  \begin{itemize}
    \item Производные вычисляются при фиксированных $z_i$.
    \item $b$ --- фиксированный скаляр, независимый от текущих сэмплов.
    \item $\E_q[\nabla\log q]=0$: такой baseline не меняет ожидание оценки.
  \end{itemize}
  \begin{block}{Почему это удобная базовая реализация}
    Нужен дифференцируемый \texttt{log\_prob}~--- производная через выбор
    компоненты и \texttt{rsample} не требуются.
  \end{block}
  \codefile{src/remix/estimators.py}
\end{frame}


\begin{frame}[fragile]{Как нужный градиент получает autograd}
  \begin{lstlisting}
z = q.sample(torch.Size([self.n_samples])).detach()
log_q = q.log_prob(z)
values = objective.evaluate(q, z)

surrogate = (values.detach() - self.baseline) * log_q + values
return EstimateResult(values.detach().mean(), -surrogate.mean())
  \end{lstlisting}
  \[
    S(\theta)=\frac1N\sum_i
       \left[(\sg(f_\theta(z_i))-b)\log q_\theta(z_i)+f_\theta(z_i)\right],
    \qquad \nabla_\theta S=\widehat g.
  \]
  \begin{itemize}
    \item \texttt{detach()} у веса исключает лишнее слагаемое производной.
    \item Минус перед средним переводит максимизацию в минимизацию.
  \end{itemize}
  \codefile{src/remix/estimators.py}
\end{frame}

\begin{frame}[fragile]{Пользовательский цикл обучения}
  \begin{lstlisting}
components = [Normal(-2.0, 1.0), Laplace(2.0, 1.0)]
logits = torch.nn.Parameter(torch.zeros(2))
target = Mixture(torch.tensor([0.2, 0.8]).log(), components)
objective = ELBOObjective(target.log_prob)
estimator = ScoreFunction(n_samples=512)
optimizer = torch.optim.Adam([logits], lr=0.05)

for step in range(300):
    optimizer.zero_grad()
    q = Mixture(logits, components)
    result = estimator.estimate(q, objective)
    result.surrogate_loss.backward()
    optimizer.step()
  \end{lstlisting}
  \medskip
  Здесь обучаются только веса~--- параметры компонент фиксированы.
  \codefile{examples/benchmark.py}
\end{frame}

\begin{frame}{Бенчмарк}
  \begin{center}
    \begin{tabular}{lll}
      \toprule
      Сценарий & Компонента 1 & Компонента 2 \\
      \midrule
      Гауссовская смесь & $\mathcal N(-2,1)$ & $\mathcal N(2,1)$ \\
      Смесь Лапласа & $\operatorname{Laplace}(-2,1)$ & $\operatorname{Laplace}(2,1)$ \\
      Разные семейства & $\mathcal N(-2,1)$ & $\operatorname{Laplace}(2,1)$ \\
      \bottomrule
    \end{tabular}
  \end{center}
  \begin{itemize}
    \item Начальные веса $(0.5,0.5)$, целевые веса $(0.2,0.8)$.
    \item Adam: шаг $0.05$, 300 итераций, 512 сэмплов на итерацию.
    \item Seed: 42; оценка KL по 10\,000 сэмплам до и после обучения.
  \end{itemize}
  \[
    \widehat{\KL}(q\Vert p)
      =\frac1M\sum_{j=1}^M[\log q(z_j)-\log p(z_j)],\quad z_j\sim q.
  \]
\end{frame}


\begin{frame}{Вывод}
  \begin{block}{Результат текущего этапа}
    Написан минимальный прототип: смесь $\to$ ELBO $\to$ score-function $\to$ стандартный оптимизатор PyTorch.
  \end{block}
  
  \begin{block}{Дальнейшие планы}
  Далее будем добавлять следующие оценщики через существующий интерфейс.
  \end{block}
\end{frame}

\end{document}
